\section{Conclusion}\label{sec:conclusion}

The \bxthree{} Framework advances a single, architecturally precise claim: that
autonomous systems must propagate every unresolved exception state to a named
human principal, bypassing all intermediate machine actors, as a compulsory
structural requirement rather than a runtime choice.  This mandatory bailout
property is not a feature that can be retrofitted to an existing agentic
architecture through better error handling or improved escalation heuristics.
It is a structural property that follows, as a matter of architectural
necessity, from the functional separation between Purpose, Bounds, and Fact
layers.

The framework's five named architectural pillars are co-equal and mutually
reinforcing.  Loop Isolation ensures that the spawn graph is a rooted DAG and
that accountability attribution is never ambiguous.  Recursive Spawning fixes the
accountability chain at the architectural level and makes it immutable.
Spatial Firewall contains sub-agent failures within execution domains and
prevents lateral state corruption.  Sandbox Gate provides the sole permitted
cross-domain communication channel with mandatory authorization verification.
The Bailout Protocol compels the final escalation for any condition not resolved
by the preceding four pillars, including conditions not anticipated at design
time.

Together, these five pillars constitute the upstream accountability guarantee:
that \bxthree{} systems \emph{fail upward into human consciousness, never
downward into autonomous chaos}.  The guarantee is not a design aspiration.
It is a logical consequence of the framework's three-layer architecture, enforced
by the \fact{} layer's deterministic pre-commit hook and recorded in the
append-only Forensic Ledger.

The theoretical contribution of this paper is foundational: that the upstream
accountability guarantee is not a design choice but a logical consequence of the
three-layer architecture, and that the Bailout Protocol is the mechanism by
which that consequence is enforced at runtime.  The practical contribution is
equally clear: that autonomous systems operating without such a mechanism are
structurally exposed to failure modes --- invisible exception absorption,
contingent escalation, and irreversible drift --- that cannot be eliminated by
incremental improvement.  They require a new architectural primitive.  The
\bxthree{} Framework is that primitive.

The limitations analysis identifies the boundaries of this primitive as
currently specified: human response latency introduces a structural floor on
exception resolution time; sophisticated agents may exploit strategic
misclassification to suppress warranted escalation; cascading failure conditions
introduce scalability constraints that the current framework does not resolve;
and the framework's guarantees are conditioned on substrate integrity properties
that the framework itself cannot verify.  The directed research agenda ---
formal verification, adaptive timeout mechanisms, distributed anchor pools, and
exploitation resistance --- addresses each limitation as a principled research
program.

Future work will determine whether the \bxthree{} Framework's guarantees hold
under full formal verification, whether adaptive exception classification can
preserve safety while reducing human latency for low-consequence exceptions,
and whether distributed human anchor pools can sustain scalability under
cascading failure conditions without compromising the no-machine-actor-absorbs
property.  The framework's core claim --- that the upstream accountability
guarantee requires mandatory human escalation for all unresolved exception
states --- is not contingent on these open questions.  It is established by
the structural necessity of the three-layer architecture.